%! Polynomial Functions and Rates of Change — Unit 2, Lesson 2.2 — Guided Notes (Answer Key)
\documentclass[10pt]{article}
\usepackage{precalculus-article}
\usepackage{precalculus-key}   % pulls in -boxes; do NOT also load -boxes
\newcommand{\TallMath}[1]{$\displaystyle #1\rule[-1.4em]{0pt}{3.2em}$}

\begin{document}

\pageheader{Unit 2, Lesson 2.2}{Guided Notes --- Answer Key}

\vspace{0.12in}

\begin{objectivebox}
By the end of this lesson, I will be able to\ldots
\begin{enumerate}[itemsep=2pt]
  \item Read the intervals where a polynomial is \ans{increasing} or
        \ans{decreasing} from a graph or a table.
  \item Tell a \ans{local} (relative) maximum or minimum from a
        \ans{global} (absolute) one.
  \item Locate a point of \ans{inflection} --- where the \textit{rate of change}
        turns around.
\end{enumerate}
\end{objectivebox}

\vspace{0.1in}

\boxguard
\begin{vocabbox}
Fill in each term as we define it together.
\par\vspace{2pt}
\noindent\textbf{\textcolor{plum}{Increasing / decreasing interval:}}\\[1pt]
\ansline{An interval where the output rises as the input rises (increasing), or falls as the input rises (decreasing).}\\[3pt]
\noindent\textbf{\textcolor{plum}{Local (relative) maximum / minimum:}}\\[1pt]
\ansline{An output value at a switch between increasing and decreasing --- or at an included endpoint of the domain.}\\[3pt]
\noindent\textbf{\textcolor{plum}{Global (absolute) maximum / minimum:}}\\[1pt]
\ansline{The greatest of all local maxima; the least of all local minima --- the extreme output over the whole domain.}\\[3pt]
\noindent\textbf{\textcolor{plum}{Concavity:}}\\[1pt]
\ansline{Concave up: the AROCs are increasing. Concave down: the AROCs are decreasing.}\\[3pt]
\noindent\textbf{\textcolor{plum}{Point of inflection:}}\\[1pt]
\ansline{An input where the rate of change switches between increasing and decreasing --- concavity changes there.}
\end{vocabbox}

\vspace{0.1in}

\boxguard
\begin{hookbox}
A trail guide lists two summits on the Ridge Trail: Bell Knob at 1{,}200 feet
and Signal Peak at 1{,}850 feet. Both are called ``summits'' --- but only one
is the highest point of the whole trail. What should we call each one?

\ansline{Each is a local (relative) maximum --- highest compared with the ground nearby.}
\ansline{Signal Peak, at 1{,}850 ft, is also the global (absolute) maximum of the trail.}
\end{hookbox}

\vspace{0.1in}

\boxguard
\begin{notesbox}{1. Polynomial Functions --- The Definition, Precisely}
A \textbf{nonconstant polynomial function} is any function that can be written
\[p(x) = a_n x^n + a_{n-1}x^{n-1} + \cdots + a_1 x + a_0,\]
where $n$ is a positive integer, every coefficient $a_i$ is a real number, and
$a_n \ne 0$.

\begin{itemize}[itemsep=4pt]
  \item The \textbf{degree} is \ans{$n$} --- the greatest exponent on $x$.
  \item The \textbf{leading term} is \ans{$a_n x^n$} and the
        \textbf{leading coefficient} is \ans{$a_n$}.
  \item A constant function such as $p(x) = 7$ is a polynomial of degree
        \ans{zero}.
\end{itemize}

\textbf{Quick check.} $f(x) = 4x^3 - x^2 + 9$: \quad
degree \ans{3} , \quad leading term \ans{$4x^3$} , \quad
leading coefficient \ans{4}
\end{notesbox}

\vspace{0.1in}

\boxguard[34]
\begin{notesbox}{2. Increasing, Decreasing, and Local Extrema}
The graph of a polynomial $f$ on the restricted domain $[0, 6]$ is below. Four
points are marked for you.

\begin{center}
\begin{tikzpicture}[x=0.78cm, y=0.5cm]
  \draw[linegray, very thin, step=1] (0,0) grid (6,7);
  \draw[->, thick] (0,0) -- (6.6,0) node[below right] {\small $x$};
  \draw[->, thick] (0,0) -- (0,7.6) node[above] {\small $f(x)$};
  \foreach \x in {1,2,3,4,5,6} \node[below] at (\x,0) {\small \x};
  \foreach \y in {1,2,3,4,5,6,7} \node[left] at (0,\y) {\small \y};
  \draw[very thick, plum] plot [smooth] coordinates
    {(0,1) (0.8,3.3) (1.5,5.2) (2,6) (2.5,5.2) (3,4.0)
     (3.5,2.6) (4,2) (4.5,2.6) (5,3.6) (5.5,5.1) (6,7)};
  \fill[goldacc] (0,1) circle (2.7pt); \node[below right] at (0.05,1) {\small $A$};
  \fill[goldacc] (2,6) circle (2.7pt); \node[above] at (2,6.15) {\small $B$};
  \fill[goldacc] (4,2) circle (2.7pt); \node[below] at (4,1.9) {\small $C$};
  \fill[goldacc] (6,7) circle (2.7pt); \node[left] at (5.85,6.9) {\small $D$};
\end{tikzpicture}
\end{center}

\begin{itemize}[itemsep=4pt]
  \item $f$ is \textbf{increasing} on \ans{$(0,2)$} and on \ans{$(4,6)$} ,
        and \textbf{decreasing} on \ans{$(2,4)$} .
  \item At $B$ the function switches from increasing to decreasing, so $B$ is
        a local \ans{maximum} . Its value is \ans{6} , occurring at
        $x = $ \ans{2} .
  \item At $C$ the function switches from decreasing to increasing, so $C$ is
        a local \ans{minimum} , with value \ans{2} at $x = $
        \ans{4} .
\end{itemize}

\textbf{Included endpoints count too.} The domain here stops at $x = 0$ and
$x = 6$, and both endpoints are included. Just to the right of $A$ the graph
rises, so $A$ is a local \ans{minimum} ; just to the left of $D$ the graph
rises into it, so $D$ is a local \ans{maximum} .

\textbf{Say it carefully:} a maximum or minimum is an \textit{output value};
it \textit{occurs at} an input value.
\end{notesbox}

\vspace{0.1in}

\boxguard
\begin{notesbox}{3. Local vs.\ Global Extrema}
Look back at the graph in Section 2 and list every local extremum you found.

\begin{itemize}[itemsep=4pt]
  \item Local maximum values: \ans{7} (at $D$) and \ans{6}
        (at $B$). The \textbf{greatest} of them is the
        \textbf{global (absolute) maximum}: \ans{7} , at $x = $
        \ans{6} .
  \item Local minimum values: \ans{1} (at $A$) and \ans{2}
        (at $C$). The \textbf{least} of them is the
        \textbf{global (absolute) minimum}: \ans{1} , at $x = $
        \ans{0} .
\end{itemize}

\textbf{The trap:} $B$ is a local maximum but \textit{not} the global one ---
the graph climbs higher later. ``Local'' means highest \textit{nearby};
``global'' means highest \textit{anywhere on the domain}.
\end{notesbox}

\vspace{0.1in}

\boxguard[30]
\begin{notesbox}{4. What the Zeros and the Degree Guarantee}
\textbf{Between two zeros, the graph has to turn around.} Below, $p$ has zeros
at $x = 1$, $x = 3$, and $x = 5$.

\begin{center}
\begin{tikzpicture}[x=0.78cm, y=0.42cm]
  \draw[linegray, very thin, step=1] (0,-3) grid (6,3);
  \draw[->, thick] (0,0) -- (6.6,0) node[below right] {\small $x$};
  \draw[->, thick] (0,-3.4) -- (0,3.6) node[above] {\small $p(x)$};
  \foreach \x in {1,2,3,4,5,6} \node[below left] at (\x,0) {\small \x};
  \draw[very thick, plum] plot [smooth] coordinates
    {(0.35,-2.7) (1,0) (1.5,1.5) (2,2) (2.5,1.5) (3,0)
     (3.5,-1.5) (4,-2) (4.5,-1.5) (5,0) (5.65,2.7)};
  \fill[plum] (1,0) circle (2.7pt);
  \fill[plum] (3,0) circle (2.7pt);
  \fill[plum] (5,0) circle (2.7pt);
\end{tikzpicture}
\end{center}

\begin{itemize}[itemsep=4pt]
  \item Between the zeros $x = 1$ and $x = 3$ there is a local
        \ans{maximum} at $x = $ \ans{2} .
  \item Between the zeros $x = 3$ and $x = 5$ there is a local
        \ans{minimum} at $x = $ \ans{4} .
  \item \textbf{Why it must happen:} the graph leaves the $x$-axis, so it
        must come back to reach the next zero --- which means it has to
        \ans{turn around} somewhere in between.
\end{itemize}

\textbf{The rule:} between every two distinct real zeros of a nonconstant
polynomial there is at least \ans{one} local maximum or minimum. So $k$
distinct real zeros force at least \ans{$k-1$} local extrema.

\medskip
\textbf{Even degree guarantees a global extremum.} Over \textit{all} real
numbers:

\smallskip
\renewcommand{\arraystretch}{1.4}
\begin{tabular}{|l|c|}
\hline
\textbf{Degree of the polynomial} & \textbf{Global maximum or global minimum?} \\
\hline
Even (like $x^2$ or $x^4 - 8x^2$) & has a global \ans{max or a global min} \\
\hline
Odd (like $x^3$ or $x^5 - x$) & \ans{has neither} \\
\hline
\end{tabular}
\smallskip

An odd-degree polynomial runs off to $+\infty$ on one side and $-\infty$ on
the other, so no output is ever the greatest or the least.
\end{notesbox}

\vspace{0.1in}

\boxguard[30]
\begin{notesbox}{5. Points of Inflection --- Where the \textit{Rate} Turns Around}
In Lesson 2.1 we watched the AROCs over consecutive equal-length intervals:
increasing AROCs mean concave up, decreasing AROCs mean concave down. Now we
name the place where that switches.

\smallskip
\renewcommand{\arraystretch}{1.4}
\begin{tabular}{|c|c|c|c|c|c|c|}
\hline
$x$ & 0 & 1 & 2 & 3 & 4 & 5 \\
\hline
$f(x)$ & 1 & 3 & 8 & 17 & 25 & 29 \\
\hline
\end{tabular}
\smallskip

AROCs over $[0,1], [1,2], [2,3], [3,4], [4,5]$: \quad
\ans{2} , \ans{5} , \ans{9} , \ans{8} , \ans{4}

\begin{itemize}[itemsep=4pt]
  \item The AROCs increase, then start to \ans{decrease} . The turn happens
        near $x = $ \ans{3} .
  \item So $f$ is concave \ans{up} before that input and concave
        \ans{down} after it. That input is a
        \textbf{point of inflection}.
\end{itemize}

\begin{center}
\begin{tikzpicture}[x=0.95cm, y=0.14cm]
  \draw[linegray, very thin, xstep=1, ystep=5] (0,0) grid (5,30);
  \draw[->, thick] (0,0) -- (5.5,0) node[below right] {\small $x$};
  \draw[->, thick] (0,0) -- (0,32.5) node[above] {\small $f(x)$};
  \foreach \x in {1,2,3,4,5} \node[below] at (\x,0) {\small \x};
  \foreach \y in {10,20,30} \node[left] at (0,\y) {\small \y};
  \draw[very thick, plum] plot [smooth] coordinates
    {(0,1) (1,3) (2,8) (3,17) (4,25) (5,29)};
  \fill[goldacc] (3,17) circle (2.7pt);
  \node[right] at (3.15,13) {\small point of inflection};
\end{tikzpicture}
\end{center}

\textbf{The contrast to keep straight:}
\begin{itemize}[itemsep=4pt]
  \item A local maximum or minimum is where the \textbf{function} turns
        around --- the AROCs change \ans{sign} .
  \item A point of inflection is where the \textbf{rate of change} turns
        around --- the AROCs stop increasing and start decreasing (or the
        reverse), but they need not change sign at all.
  \item In the table above, every AROC is positive, so $f$ is increasing the
        whole way and has \ans{no} local extrema --- yet it still has a
        point of inflection.
\end{itemize}
\end{notesbox}

\vspace{0.1in}

\boxguard
\begin{practicebox}
\begin{enumerate}[itemsep=8pt]
  \item A polynomial $g$ has exactly four distinct real zeros.

        \begin{enumerate}[label=(\alph*), itemsep=4pt]
          \item At least how many local extrema must $g$ have?
                \ans{3}
          \item Explain in one sentence.

                \vspace{0.05in}
                \ansline{Each of the three gaps between consecutive zeros forces at least one local max or min.}
        \end{enumerate}

  \item A polynomial $h$ has AROCs $4,\ 9,\ 15,\ 11,\ 6$ over five
        consecutive equal-length intervals covering $[0, 5]$.

        \begin{enumerate}[label=(\alph*), itemsep=4pt]
          \item Does $h$ have a local maximum or minimum on $[0,5]$?
                \ans{No --- all AROCs are positive}
          \item Near which input is there a point of inflection?
                $x = $ \ans{3}
        \end{enumerate}
\end{enumerate}
\end{practicebox}

\end{document}
